\section{Print, Trent, and the Roman Editions}

Movable-type printing changed the scale and economics of biblical reproduction. It did not make the first printed exemplar critically ancient. The Gutenberg Bible, completed probably late in 1455 at Mainz, printed a Latin Vulgate in the textual world of late-medieval manuscripts. Its achievement was technological and material: many substantially identical copies could now descend from one typeset text, while rubrication and decoration might still be completed by hand (Library of Congress, ``The Gutenberg Bible'').

\subsection{Print fixes a copy and multiplies its choices}

Early printers selected from available manuscripts and earlier printings. Once typeset, a reading that had been local could be multiplied widely; compositor errors could be multiplied with it. Corrected reprints, humanist comparison with Greek and Hebrew, and new Latin translations consequently proliferated together. Erasmus's work on the Greek New Testament and his revised Latin, the Complutensian Polyglot, and other sixteenth-century editions made discrepancies among source languages and Latin forms newly visible (Houghton 2016, 108--110).

The age of print therefore sharpened two different questions. Scholars asked which readings had the best historical support. Churches and schools also needed a stable Latin text for public use. A single edition might serve both needs imperfectly, but the needs are not identical.

\subsection{What Trent declared in 1546}

On 8 April 1546 the Council of Trent's fourth session issued a decree on the edition and use of the sacred books. Among the Latin editions then circulating, it declared the old and Vulgate edition---approved by long ecclesial use---to be held as authentic in public readings, disputations, preaching, and exposition. The decree also required Scripture, especially the old Vulgate, to be printed as accurately as possible (Council of Trent, \emph{Decretum de editione et usu sacrorum librorum}).

The syntax and setting matter. The comparison is among Latin editions; the uses are named; and the demand for a more accurate printing acknowledges that no single flawless printed text had simply been identified. The decree did not declare every transmitted Vulgate reading inspired, make Latin superior to the Hebrew and Greek sources, or identify the Clementine edition decades before it existed. In 1943 Pius XII expressly explained the Vulgate's Tridentine authenticity as juridical rather than primarily critical and affirmed recourse to the original-language texts (\emph{Divino afflante Spiritu}, 20--22).

\evidenceframe{Council of Trent, session IV, \emph{Decretum de editione et usu sacrorum librorum}, 8 April 1546.}{The old Vulgate was to be held authentic among circulating Latin editions for specified public ecclesial uses, and Scripture was to be printed as accurately as possible.}{The decree does not select a fourth-century manuscript, name the later Clementine text, abolish Hebrew and Greek, or make a claim about every philological detail.}

The decree on the scriptural canon, issued in the same session, is a separate act. It identifies books and parts received by the Catholic Church; the edition-and-use decree governs the status of a Latin edition. Canonical reception and textual reconstruction must not be made synonyms.

\subsection{From commission work to the Sixtine edition}

Roman work on a corrected Vulgate continued under several popes and drew on manuscripts, printed editions, source-language texts, and patristic evidence. Pope Sixtus V issued the first papally promulgated printed Vulgate in 1590. It is known as the Sixtine or Sistine edition. Sixtus died that year, and his edition did not become the durable settlement.

The later Clementine preface narrates the transition diplomatically. It says Sixtus had judged that printing faults required renewed work, which death prevented him from completing; work continued through the short pontificates of Gregory XIV and Innocent IX and was completed under Clement VIII. Modern histories also note the recall of the 1590 printing and debate the mix of textual, political, and reputational motives. The secure publication sequence is more important here than a single motive: Sixtine 1590, then a materially different Roman edition under Clement VIII in 1592 (Clementine preface; Houghton 2016, 132--133).

\subsection{The Clementine settlement}

The 1592 edition, followed by corrected Roman printings in 1593 and 1598, is called Clementine or Sixto-Clementine because its title continued to honor Sixtus's command. Its anonymous preface describes the received Bible candidly as partly Jerome's translation or correction and partly retained from an older Latin edition. It also says the work was corrected with as much diligence as possible. These statements undercut two later myths at once: sole Hieronymian authorship and a claim that editorial perfection was easy to guarantee.

The Clementine text gave Catholic public use a stable printed norm for centuries. Stability did not make it a diplomatic transcript of any medieval codex or a critical reconstruction of Jerome's earliest attainable text. Its readings arose from an early-modern editorial settlement, and its book order, chapter and verse system, punctuation, orthography, prologues, and appendix belong to that edition's own history. The Prayer of Manasseh and the books called 3 and 4 Esdras were printed in an appendix after the New Testament so that they would not disappear entirely; that placement did not add them to Trent's canonical list.

\claimcheck{``Trent declared the Clementine Vulgate perfect''}{Trent acted in 1546, forty-six years before the Clementine edition. It conferred juridical authenticity on the old Vulgate among Latin editions in specified public uses and called for accurate printing. The Clementine text was a later Roman editorial answer, not a text named or declared verbally perfect by the council.}

The Clementine's long use makes it a major historical object in its own right. It is often closer than an earliest-text critical edition to the wording encountered by early-modern and later Catholic writers. A historian quoting a seventeenth-century theologian may therefore need the Clementine; an editor reconstructing a fifth-century Vulgate book needs a different instrument. Calling both tasks ``using the Vulgate'' without naming an edition obscures the evidence.
